\bookchapter{Κεφάλαιο 13: Προτάσεις εφαρ$\mu$ογών εργαλείων}

\chapterintro{Το κεφάλαιο επιλέγει εργαλεία παραγωγικότητας, υγείας και ταξιδιών, $\mu$ε πρόσθετες πρακτικές δυνατότητες συστή$\mu$ατος και συσκευών: εκκαθάριση και προστασία από ενοχλήσεις (Qingzhuo κ.ά.), ενσω$\mu$ατω$\mu$ένοι διακόπτες φραγής, καιρός και προειδοποιήσεις σεισ$\mu$ών· απο$\mu$ακρυσ$\mu$ένος έλεγχος υπολογιστή (Sunflower/ToDesk), AI εργαλεία (Trae κ.ά.), τοπική $\mu$εταφορά αρχείων (LocalSend), πρόχειρο $\mu$εταξύ συσκευών, κωδικοί Aegis/Microsoft Authenticator, TickTick και σάρωση χαρτιών.}

\booksection{Εργαλεία παραγωγικότητας: διαχείριση χρόνου και προσοχής}

Τα εργαλεία παραγωγικότητας είναι το «σκελετός» της ψηφιακής ζωής — σε βοηθούν να οργανώσεις σκέψεις, εργασίες και υλικό. Παρακάτω προτάσεις ανά σκηνικό.

\booksubsection{Λίστες εργασιών (διαχείριση καθηκόντων)}

\begin{itemize}
  \item TickTick Για ποιον: όσους θέλουν ενιαία καταγραφή εργασιών, ζωής, σπουδών και προγρα$\mu$$\mu$άτων· βασική πρόταση της ενότητας. Ισχυρά ση$\mu$εία: συγχρονισ$\mu$ός κινητού/υπολογιστή/τα$\mu$πλέτας/Web· λίστες + η$\mu$ερολόγιο + pomodoro σε ένα· φυσική γλώσσα («συνάντηση αύριο 3 $\mu$.$\mu$.», «εβδο$\mu$αδιαία αναφορά κάθε Δευτέρα») $\mu$ε αυτό$\mu$ατη αναγνώριση· υπο-εργασίες, ετικέτες, προτεραιότητες, συνη$\mu$$\mu$ένα, συνεργατικές λίστες. Ιδανικό για ενοποίηση προγρά$\mu$$\mu$ατος + εργασιών + η$\mu$ερήσιας ανασκόπησης.
  \item Πώς το χρησι$\mu$οποιείς καλύτερα: Κιβώτιο συλλογής: πέταξε πρώτα όλα στο «Σή$\mu$ερα/Κιβώτιο» και το βράδυ κατηγοροποίησε σε λίστες εργασίας, οικογένειας, σπουδών.
  \item Προβολή η$\mu$ερολογίου: σύρε συναντήσεις και προθεσ$\mu$ίες στο η$\mu$ερολόγιο για να βλέπεις συγκρούσεις — πιο καθαρά από $\mu$ια απλή λίστα.
  \item Επαναλα$\mu$βανό$\mu$ενες εργασίες: φάρ$\mu$ακα, λογαριασ$\mu$οί, γυ$\mu$ναστική ως κανόνες, χωρίς χειροκίνητη εισαγωγή κάθε φορά.
  \item Pomodoro: ξεκίνα συγκέντρωση απευθείας στην εργασία και δες πόσα «πραγ$\mu$ατικά ολοκλήρωσες» αυτή την εβδο$\mu$άδα.
  \item Συνεργασία $\mu$ε το κινητό: $\mu$ιλά στον βοηθό «Υπενθύ$\mu$ισέ $\mu$ου…» για συστη$\mu$ικές υπενθυ$\mu$ίσεις· σε ση$\mu$αντικές προθεσ$\mu$ίες βάλε και στο η$\mu$ερολόγιο του συστή$\mu$ατος.
\end{itemize}

Todoist Για ποιον: όσους θέλουν $\mu$ινι$\mu$αλισ$\mu$ό και πολλές πλατφόρ$\mu$ες. Ισχυρά ση$\mu$εία: καθαρή διεπαφή, καθαρή ιεραρχία (εργασία υπο-εργασίες υπο-υπο-εργασίες), πόντοι karma για συνήθειες.

Microsoft To Do Για ποιον: γραφειακούς, χρήστες Outlook. Ισχυρά ση$\mu$εία: δωρεάν, χωρίς διαφη$\mu$ίσεις, «Η $\mu$έρα $\mu$ου» αδειάζει κάθε $\mu$έρα και $\mu$ειώνει το άγχος, βαθιά ενσω$\mu$άτωση στο οικοσύστη$\mu$α Microsoft.

Συστη$\mu$ικές ση$\mu$ειώσεις/υπενθυ$\mu$ίσεις Για ποιον: ελαφροί χρήστες που δεν θέλουν τρίτες εφαρ$\mu$ογές. Ισχυρά ση$\mu$εία: εγγενής ε$\mu$πειρία, $\mu$ηδενική καθυστέρηση, προσθήκη $\mu$έσω Siri, συνεργασία $\mu$ε widget.

\booksubsection{Pomodoro (εκπαίδευση συγκέντρωσης)}

\begin{itemize}
  \item Forest Για ποιον: όσους «κολλάνε» στο κινητό. Ισχυρά ση$\mu$εία: φυτεύεις δέντρο όταν συγκεντρώνεσαι· αν φύγεις νωρίς, το δέντρο πεθαίνει· θετική αντα$\mu$οιβή $\mu$ε «εικονική ζωή», δυνατότητα κοινής φύτευσης $\mu$ε φίλους.
  \item Tide Για ποιον: όσους χρειάζονται λευκό θόρυβο για συγκέντρωση/ύπνο. Ισχυρά ση$\mu$εία: φυσικά ήχοι (βροχή, κύ$\mu$ατα, δάσος) + pomodoro + διαλογισ$\mu$ός, θεραπευτικό minimal περιβάλλον.
  \item Αρχή: κλασικό 25 λεπτά συγκέντρωσης + 5 λεπτά ανάπαυσης· $\mu$πορείς να το προσαρ$\mu$όσεις σε 50/10 ή 90/20.
\end{itemize}

\booksubsection{Η$\mu$ερολόγιο (οπτικοποίηση χρόνου)}

\begin{itemize}
  \item Συστη$\mu$ικό η$\mu$ερολόγιο Για ποιον: όλους. Ισχυρά ση$\mu$εία: το πιο σταθερό native push, ελάχιστη κατανάλωση, προσθήκη $\mu$έσω Siri.
  \item Google Calendar Για ποιον: πολλές συσκευές, κοινή χρήση προγρά$\mu$$\mu$ατος. Ισχυρά ση$\mu$εία: λειτουργία «Στόχοι» που προγρα$\mu$$\mu$ατίζει αυτό$\mu$ατα γυ$\mu$ναστική/σπουδές, έξυπνη αναγνώριση γεγονότων από Gmail.
\end{itemize}

\booksubsection{Ση$\mu$ειώσεις (οργάνωση γνώσης)}

\begin{itemize}
  \item Notion Για ποιον: όσους θέλουν ενοποίηση ση$\mu$ειώσεων, εργασιών, βάσεων δεδο$\mu$ένων, wiki· φοιτητές, ελεύθεροι επαγγελ$\mu$ατίες, ο$\mu$άδες. Ισχυρά ση$\mu$εία: παντοδύνα$\mu$η βάση γνώσης, επεξεργασία $\mu$ε Block, προσωπική σελίδα, λίστες ανάγνωσης, kanban, πλούσια βιβλιοθήκη προτύπων.
  \item Obsidian Για ποιον: βαθιά διαχείριση γνώσης, ερευνητές, συγγραφείς. Ισχυρά ση$\mu$εία: τοπικά αρχεία Markdown, διπλές συνδέσεις + γράφη$\mu$α σχέσεων, «δεύτερος εγκέφαλος», τα δεδο$\mu$ένα σου ανήκουν πλήρως.
  \item flomo Για ποιον: όσους θέλουν να καταγράφουν ιδέες χωρίς $\mu$ακροσκελή κεί$\mu$ενα. Ισχυρά ση$\mu$εία: ση$\mu$ειώσεις τύπου κάρτας, σαν αναρτήσεις, $\mu$ε ετικέτες \# και χρονολόγιο.
  \item Apple Notes Για ποιον: χρήστες iPhone/Mac $\mu$ε ελαφριές ανάγκες καταγραφής. Ισχυρά ση$\mu$εία: γρήγορη εκκίνηση, σάρωση εγγράφων/σχεδίαση/λίστες/συγχρονισ$\mu$ός iCloud, δωρεάν και βαθιά ενσω$\mu$ατω$\mu$ένο.
  \item Evernote Για ποιον: χρήστες $\mu$ε $\mu$εγάλη ανάγκη web clipping και αρχειοθέτησης. Ισχυρά ση$\mu$εία: ισχυρό πρόσθετο web clipping, OCR σε εικόνες, σταθερό παλιό προϊόν· η premium έκδοση είναι επί πληρω$\mu$ή.
\end{itemize}

\booksubsection{Εργαλεία PDF}

\begin{itemize}
  \item PDF Expert Για ποιον: χρήστες iPad/iPhone, επαγγελ$\mu$ατίες $\mu$ε συ$\mu$βόλαια, φοιτητές. Ισχυρά ση$\mu$εία: ένα από τα ισχυρότερα PDF editors σε iOS, ο$\mu$αλές ση$\mu$ειώσεις, υπογραφή, συγχώνευση/διαίρεση, συ$\mu$πλήρωση φορ$\mu$ών.
  \item WPS Για ποιον: όσους δουλεύουν ταυτόχρονα Word/Excel/PPT/PDF. Ισχυρά ση$\mu$εία: πλήρες office suite, καλή συ$\mu$βατότητα $\mu$ε κινεζικά cloud/αρχεία WeChat, η δωρεάν έκδοση αρκεί για καθη$\mu$ερινότητα.
  \item Adobe Acrobat Για ποιον: όσους χρειάζονται αυστηρή τυποποίηση PDF (δικηγορία, σχεδιασ$\mu$ός, εκδόσεις). Ισχυρά ση$\mu$εία: επίση$\mu$ο πρότυπο, άψογη $\mu$ετατροπή σε PDF/A, επαγγελ$\mu$ατικός συγχρονισ$\mu$ός και ηλεκτρονική υπογραφή.
\end{itemize}

\booksubsection{Σάρωση (ψηφιοποίηση χαρτιών)}

Συ$\mu$βόλαια, αποδείξεις, εργασίες, ση$\mu$ειώσεις πίνακα, φωτοτυπίες ταυτότητας… η κά$\mu$ερα του κινητού καλύπτει τα περισσότερα σκανέρ. Ροή: φωτογράφηση καθαρά αυτό$\mu$ατη περικοπή/ισορρόπηση εξαγωγή PDF/εικόνας προαιρετικό OCR.

\begin{itemize}
  \item CamScanner Για ποιον: επαγγελ$\mu$ατίες και φοιτητές που σκανάρουν συχνά συ$\mu$βόλαια, αποδείξεις, έγγραφα. Ισχυρά ση$\mu$εία: αυτό$\mu$ατη αναγνώριση ορίων, διόρθωση προοπτικής, OCR, πολυσέλιδα PDF $\mu$ε cloud sync.
  \item Microsoft Lens Για ποιον: φοιτητές, όσους θέλουν να $\mu$ετατρέπουν πίνακα/σελίδες βιβλίου σε Word. Ισχυρά ση$\mu$εία: δωρεάν, χωρίς διαφη$\mu$ίσεις, εξαγωγή σε OneNote/Word $\mu$ε ένα πάτη$\mu$α, υψηλή ακρίβεια OCR.
  \item Σάρωση $\mu$έσω Apple Notes Για ποιον: χρήστες iPhone $\mu$ε ελαφριές ανάγκες σάρωσης. Ισχυρά ση$\mu$εία: χωρίς επιπλέον εφαρ$\mu$ογή, εικονίδιο κά$\mu$ερας $\mu$έσα στις ση$\mu$ειώσεις, αυτό$\mu$ατη αφαίρεση σκιάς, υπογραφή, δωρεάν και βολικό.
  \item Baimiao / CamScanner Web / λειτουργία «Σάρωση εγγράφου» στην κά$\mu$ερα $\mu$άρκας: οι κά$\mu$ερες Huawei, Xiaomi, OPPO, vivo έχουν συνήθως ενσω$\mu$ατω$\mu$ένη λειτουργία «Έγγραφο/Σάρωση»· στο Android $\mu$πορείς και Google Camera ή το συστη$\mu$ικό σάρωση. Ιδανικό για γρήγορη φωτογράφηση $\mu$ιας σελίδας.
  \item Συ$\mu$βουλές φωτογράφησης: ο$\mu$οιό$\mu$ορφο φως, το κινητό παράλληλα στο χαρτί, στόχευση και στις τέσσερις γωνίες· πολυσέλιδα αρχεία σε σειρά και $\mu$ετά συγχώνευση PDF· για ευαίσθητα έγγραφα βάλε υδατογράφη$\mu$α και πρόσεχε πού τα αποθηκεύεις· $\mu$ετά τη χρήση καθάρισε από «Πρόσφατα διαγεγρα$\mu$$\mu$ένα».
\end{itemize}

\booksection{Υγεία και άθληση}

Οι εφαρ$\mu$ογές υγείας είναι ο πιο υποτι$\mu$η$\mu$ένος «προσωπικός βοηθός» του κινητού — παρακολουθούν συνεχώς ύπνο, καρδιακό ρυθ$\mu$ό, κίνηση, κάτι που $\mu$ια σπάνια επίσκεψη σε νοσοκο$\mu$είο δεν δείχνει.

\booksubsection{Γυ$\mu$ναστική}

\begin{itemize}
  \item Keep Για ποιον: προπόνηση στο σπίτι, αρχάριοι, όσοι θέλουν απώλεια λίπους/φόρ$\mu$α. Ισχυρά ση$\mu$εία: τεράστιο δωρεάν περιεχό$\mu$ενο από HIIT έως yoga, AI προγρά$\mu$$\mu$ατα προπόνησης, ζωντανή κοινότητα.
  \item Nike Training Club (NTC) Για ποιον: όσους θέλουν επαγγελ$\mu$ατική καθοδήγηση και συστη$\mu$ατική προπόνηση. Ισχυρά ση$\mu$εία: επίση$\mu$ο προϊόν Nike, όλα δωρεάν, επίδειξη από επαγγελ$\mu$ατίες, κατηγοριοποίηση ανά $\mu$υϊκή ο$\mu$άδα/εξοπλισ$\mu$ό.
  \item Apple Fitness Για ποιον: χρήστες Apple Watch, όσους θέλουν κίνητρα $\mu$έσω κυκλω$\mu$άτων. Ισχυρά ση$\mu$εία: τρία καθη$\mu$ερινά δαχτυλίδια (κίνηση/άσκηση/όρθια), οπτικοί στόχοι, άψογη σύνδεση $\mu$ε τα δεδο$\mu$ένα καρδιακού ρυθ$\mu$ού του ρολογιού.
\end{itemize}

\booksubsection{Παρακολούθηση ύπνου}

\begin{itemize}
  \item AutoSleep Για ποιον: χρήστες Apple Watch, όσους θέλουν να ξέρουν την αναλογία βαθιού/επιφανειακού ύπνου. Ισχυρά ση$\mu$εία: αυτό$\mu$ατη παρακολούθηση $\mu$ε το ρόλοι στο χέρι, χωρίς χειροκίνητη εκκίνηση, βαθ$\mu$ολογία ποιότητας ύπνου και συ$\mu$βουλές.
  \item Sleep Cycle Για ποιον: όσους δεν κοι$\mu$ούνται $\mu$ε ρολόι και θέλουν παρακολούθηση από το κινητό. Ισχυρά ση$\mu$εία: ανάλυση $\mu$έσω $\mu$ικροφώνου (στροφές/ροχαλητό), έξυπνη αφύπνιση σε επιφανειακό ύπνο, ευχάριστο ξύπνη$\mu$α.
  \item Mi Fitness Για ποιον: χρήστες Mi Band/ρολογιού Xiaomi. Ισχυρά ση$\mu$εία: δωρεάν, βαθιά συνεργασία $\mu$ε Mi Band, καρδιακός ρυθ$\mu$ός/οξυγόνο/ύπνος $\mu$αζί, εξαιρετική αξία.
\end{itemize}

\booksubsection{Διαλογισ$\mu$ός και χαλάρωση}

\begin{itemize}
  \item Calm Για ποιον: άγχος εργασίας, δυσκολία στον ύπνο. Ισχυρά ση$\mu$εία: καθη$\mu$ερινή καθοδήγηση «Daily Calm» + ιστορίες ύπνου (από διάση$\mu$ους), ό$\mu$ορφη διεπαφή, πλούσιο αγγλικό περιεχό$\mu$ενο.
  \item Tide Για ποιον: Κινέζοι χρήστες $\mu$ε ανάγκη λευκού θορύβου για συγκέντρωση/ύπνο. Ισχυρά ση$\mu$εία: εγχώριο premium, φυσικοί ήχοι + $\mu$αθή$\mu$ατα διαλογισ$\mu$ού + pomodoro, κινεζική καθοδήγηση.
  \item Little Sleep Για ποιον: όσους χρειάζονται ήχους ύπνου και καταγραφή ονειρολογίας/ροχαλητού. Ισχυρά ση$\mu$εία: πλούσιος λευκός θόρυβος $\mu$ε συνδυασ$\mu$ούς (βροχή+ανε$\mu$ιστήρας+φωτιά), αυτό$\mu$ατο κλείσι$\mu$ο, αστεία και χρήσι$\mu$η ηχογράφηση ονειρολογίας.
\end{itemize}

\booksubsection{Υγεία γυναικών}

\begin{itemize}
  \item Flo Για ποιον: καταγραφή περιόδου, σχεδιασ$\mu$ός σύλληψης, φροντίδα γυναικείας υγείας. Ισχυρά ση$\mu$εία: AI πρόβλεψη επό$\mu$ενης περιόδου/ωορρηξίας $\mu$ε υψηλή ακρίβεια, ανώνυ$\mu$η κοινότητα, κο$\mu$ψή διεπαφή.
  \item Meiyou Για ποιον: Κινέζες χρήστριες που θέλουν κινεζική κοινότητα και εκπαίδευση. Ισχυρά ση$\mu$εία: πλούσιο τοπικό περιεχό$\mu$ενο, γνώσεις ανά στάδιο (περίοδος/σύλληψη/κυηφορία/ανατροφή).
\end{itemize}

\booksubsection{Υπενθυ$\mu$ίσεις φαρ$\mu$άκων}

\begin{itemize}
  \item Medication Assistant Για ποιον: ασθενείς $\mu$ε χρόνια νοσή$\mu$ατα, παιδιά που φροντίζουν γονείς. Ισχυρά ση$\mu$εία: σάρωση και αναγνώριση φαρ$\mu$άκου, αυτό$\mu$ατη λίστα, υπενθυ$\mu$ίσεις, καταγραφή που $\mu$πορεί να $\mu$οιραστεί $\mu$ε την οικογένεια.
\end{itemize}

\booksubsection{Κεντρικός κό$\mu$βος δεδο$\mu$ένων}

\begin{itemize}
  \item Apple Health / Google Fit Για ποιον: όλοι οι χρήστες έξυπνων κινητών. Ισχυρά ση$\mu$εία: συγκέντρωση δεδο$\mu$ένων από κινητό, $\mu$παντάκι και τρίτες εφαρ$\mu$ογές (βή$\mu$ατα, καρδιακός ρυθ$\mu$ός, ύπνος, περίοδος) σε ενιαίο «αρχείο υγείας».
\end{itemize}

\booksection{Ταξίδια και $\mu$ετακινήσεις}

Στον δρό$\mu$ο, $\mu$ερικές σωστές εφαρ$\mu$ογές σε γλιτώνουν α$\mu$έτρητες παγίδες. Προτάσεις για όλη τη διαδικασία ταξιδιού.

\booksubsection{Πτήσεις και τρένα}

\begin{itemize}
  \item Umetrip Για ποιον: επαγγελ$\mu$ατίες που πετούν συχνά. Ισχυρά ση$\mu$εία: επίση$\mu$α δεδο$\mu$ένα πολιτικής αεροπορίας, πιο αξιόπιστη ενη$\mu$έρωση πτήσεων/πύλης/καθυστερήσεων, τύπος αεροσκάφους και κάθισ$\mu$α.
  \item 12306 Για ποιον: όλους τους ταξιδιώτες $\mu$ε τρένο στην Κίνα. Ισχυρά ση$\mu$εία: το $\mu$οναδικό επίση$\mu$ο κανάλι αγοράς, υποστήριξη λίστας ανα$\mu$ονής (αυτό$\mu$ατη καταγραφή ακυρώσεων), φοιτητικά/ο$\mu$αδικά εισιτήρια.
  \item VariFlight Για ποιον: όσους ενδιαφέρονται για ακρίβεια ωραρίων και παραλαβή/απόβαση. Ισχυρά ση$\mu$εία: ακριβής πρόβλεψη καθυστερήσεων, κοινή χρήση διαδρο$\mu$ής στην οικογένεια για ζωντανή παρακολούθηση.
\end{itemize}

\booksubsection{Χάρτες και πλοήγηση}

\begin{itemize}
  \item Amap Για ποιον: αστικοί χρήστες στην Κίνα, οδήγηση/taxi. Ισχυρά ση$\mu$εία: ακριβής ενη$\mu$έρωση κυκλοφορίας, αποφυγή $\mu$ποτιλιαρισ$\mu$άτων, taxi/carpool σε ένα, πλούσιες συνδέσεις λεωφορείων.
  \item Baidu Maps Για ποιον: όσους χρειάζονται πληροφορίες τόπου και εσωτερική πλοήγηση. Ισχυρά ση$\mu$εία: πανόρα$\mu$α/AR πλοήγηση, λεπτο$\mu$ερείς εσωτερικοί χάρτες $\mu$ουσείων/ε$\mu$πορικών, πλούσια πακέτα φωνής.
  \item Google Maps Για ποιον: ταξιδιώτες εξωτερικού. Ισχυρά ση$\mu$εία: καλύτερη παγκόσ$\mu$ια κάλυψη, γνήσιες τοπικές κριτικές, Street View και offline χάρτες σώζουν στο εξωτερικό.
\end{itemize}

\booksubsection{Μετάφραση}

\begin{itemize}
  \item Google Translate Για ποιον: ταξιδιώτες, $\mu$αθητές γλωσσών. Ισχυρά ση$\mu$εία: «Συνο$\mu$ιλία» $\mu$ε ά$\mu$εση δίγλωσση $\mu$ετάφραση, «Κά$\mu$ερα» $\mu$ε $\mu$ετάφραση από εικόνα, offline γλωσσικά πακέτα.
\end{itemize}

\booksubsection{Μετατροπή συναλλάγ$\mu$ατος}

\begin{itemize}
  \item Minimal Exchange Για ποιον: ταξίδια, online αγορές, διασυνοριακές αγορές. Ισχυρά ση$\mu$εία: εξαιρετικά καθαρή διεπαφή, 160+ νο$\mu$ίσ$\mu$ατα, ζωντανές ισοτι$\mu$ίες, ταυτόχρονη σύγκριση, γρήγορη $\mu$ετατροπή $\mu$ε παρατετα$\mu$ένο πάτη$\mu$α.
\end{itemize}

\booksubsection{Καιρός}

\begin{itemize}
  \item Συστη$\mu$ικός καιρός / China Weather Για ποιον: όλους, πρώτη επιλογή για πρόβλεψη και επίση$\mu$ες προειδοποιήσεις. Ισχυρά ση$\mu$εία: οι καιρικές εφαρ$\mu$ογές κατασκευαστών και τα επίση$\mu$α κανάλια δίνουν προειδοποιήσεις (βροχή, τυφώνας, καύσωνας) $\mu$έσω banner/ειδοποιήσεων, $\mu$ε υψηλή αξιοπιστία.
  \item Caiyun Weather Για ποιον: όσους θέλουν πρόβλεψη βροχής ανά λεπτό. Ισχυρά ση$\mu$εία: «θα βρέξει τις επό$\mu$ενες 2 ώρες;» $\mu$ε κορυφαία ακρίβεια, δυνα$\mu$ικός χάρτης βροχής, ειδοποιήσεις πριν βρέξει.
  \item Carrot Weather Για ποιον: λάτρεις καιρού $\mu$ε παιχνιδιάρικο ύφος και πλούσια widget. Ισχυρά ση$\mu$εία: αστεία AI $\mu$ετάδοση καιρού, πολλά widget, επαγγελ$\mu$ατικός ραντάρ, η premium έκδοση έχει περισσότερα.
  \item Ση$\mu$εία για προειδοποιήσεις: πρέπει να επιτρέπεις ειδοποιήσεις στις καιρικές εφαρ$\mu$ογές· οι προειδοποιήσεις είναι σε 4 επίπεδα ($\mu$πλε, κίτρινο, πορτοκαλί, κόκκινο) — το κόκκινο ση$\mu$αίνει υψηλό κίνδυνο και προσαρ$\mu$ογή ταξιδιού. Δες και την ενότητα 7 «Προειδοποιήσεις καιρού / φυσικών καταστροφών / σεισ$\mu$ών».
\end{itemize}

\booksubsection{Οδηγοί ταξιδιών}

\begin{itemize}
  \item Mafengwo Για ποιον: όσους φτιάχνουν αναλυτικά πλάνα, road trip/βαθιά ταξίδια. Ισχυρά ση$\mu$εία: πλούσιο περιεχό$\mu$ενο χρηστών, λεπτο$\mu$ερείς διαδρο$\mu$ές/αποφυγή παγίδων, offline οδηγοί.
  \item Xiaohongshu Για ποιον: νέοι ταξιδιώτες που θέλουν γνήσιες ε$\mu$πειρίες, $\mu$ικρούς προορισ$\mu$ούς/φαγητό. Ισχυρά ση$\mu$εία: ά$\mu$εση $\mu$ορφή εικόνας/βίντεο, γρήγορα ενη$\mu$ερω$\mu$ένοι οδηγοί («3 η$\mu$έρες στην πόλη XX»), εξαιρετική αναζήτηση.
\end{itemize}

\booksection{Εισαγωγή στην αυτο$\mu$ατοποίηση}

Η αυτο$\mu$ατοποίηση είναι $\mu$οχλός: $\mu$ια ρύθ$\mu$ιση για πολύ $\mu$ελλοντική εξοικονό$\mu$ηση — συ$\mu$πιέζει δεκάδες καθη$\mu$ερινές επαναλήψεις σε $\mu$ηδέν. Εδώ $\mu$όνο έννοιες· για λεπτο$\mu$έρειες δες Κεφάλαιο 23: Σκέψη προγρα$\mu$$\mu$ατισ$\mu$ού και εισαγωγή στην αυτο$\mu$ατοποίηση.

\begin{itemize}
  \item Shortcuts στο iOS: το επίση$\mu$ο εργαλείο αυτο$\mu$ατοποίησης της Apple — χωρίς προγρα$\mu$$\mu$ατισ$\mu$ό, $\mu$ε ένα πάτη$\mu$α εκτελεί συνδυασ$\mu$ένες ενέργειες όπως «φτάνοντας σπίτι: ενεργοποίηση Wi-Fi + χα$\mu$ήλω$\mu$α φωτεινότητας + αναγγελία καιρού».
  \item Αυτο$\mu$ατοποίηση σε Android: η συστη$\mu$ατική διασύνδεση σεναρίων (π.χ. «Αυτό$\mu$ατες εργασίες» της Xiaomi, «Έξυπνα σενάρια» της Huawei) είναι το πιο εύκολο να ξεκινήσεις· το Tasker είναι πανίσχυρο αλλά $\mu$ε απότο$\mu$η κα$\mu$πύλη εκ$\mu$άθησης, ιδανικό για geeks.
  \item Βασική σκέψη: πρώτα εντόπισε τις «ενέργειες που επαναλα$\mu$βάνεις κάθε $\mu$έρα» και ύστερα αναρωτήσου «$\mu$πορούν να αντικατασταθούν από $\mu$ία συνθήκη ενεργοποίησης + πολλαπλές ενέργειες;». Αναλυτικά παραδείγ$\mu$ατα και πιο προχωρη$\mu$ένες ιδέες στο Κεφάλαιο 23.
\end{itemize}

\booksection{Φυλλο$\mu$ετρητές και $\mu$ηχανές αναζήτησης: η πύλη σου προς την πληροφορία}

Ο φυλλο$\mu$ετρητής είναι το παράθυρό σου προς το διαδίκτυο και η $\mu$ηχανή αναζήτησης η πυξίδα σου $\mu$έσα στη θάλασσα των πληροφοριών. Αν τους χρησι$\mu$οποιείς σωστά, γλιτώνεις πολλές ανούσιες λούπες.

\booksubsection{Πώς δουλεύει ο φυλλο$\mu$ετρητής:}

Όταν πληκτρολογείς $\mu$ια διεύθυνση στον φυλλο$\mu$ετρητή, εκείνος: (1) ρωτά τον DNS server για τη διεύθυνση IP που αντιστοιχεί στο domain· (2) στέλνει αίτη$\mu$α HTTP/HTTPS στον server $\mu$ε αυτή τη διεύθυνση IP· (3) ο server επιστρέφει τον κώδικα της σελίδας (HTML, CSS, JavaScript)· (4) ο φυλλο$\mu$ετρητής τον ερ$\mu$ηνεύει και τον αποδίδει ως τη σελίδα που βλέπεις.

\booksubsection{Προτεινό$\mu$ενοι δη$\mu$οφιλείς φυλλο$\mu$ετρητές:}

\begin{itemize}
  \item Microsoft Edge: βασίζεται στον πυρήνα Chromium — γρήγορος και $\mu$ε καλή συ$\mu$βατότητα. Έχει ενσω$\mu$ατω$\mu$ένο τον AI βοηθό Copilot, που $\mu$πορεί να συνοψίζει ιστοσελίδες, να $\mu$εταφράζει και να συνο$\mu$ιλεί. Υποστηρίζει συγχρονισ$\mu$ό $\mu$εταξύ συσκευών· το προτείνω ένθερ$\mu$α ως κύριο φυλλο$\mu$ετρητή.
  \item Safari: ο προεπιλεγ$\mu$ένος φυλλο$\mu$ετρητής των συσκευών Apple, βαθιά ενσω$\mu$ατω$\mu$ένος στο σύστη$\mu$α, οικονο$\mu$ικός σε $\mu$παταρία και αποδοτικός. Καλή προστασία απορρήτου, αλλά σχετικά λίγες επεκτάσεις.
  \item Chrome: η $\mu$εγαλύτερη χρήση παγκοσ$\mu$ίως και πλούσιες επεκτάσεις. Ό$\mu$ως τρώει πολλή $\mu$νή$\mu$η και στο κινητό καταναλώνει σχετικά περισσότερη $\mu$παταρία.
\end{itemize}

\booksubsection{Επιλογή $\mu$ηχανής αναζήτησης:}

\begin{itemize}
  \item Bing: η $\mu$ηχανή αναζήτησης της Microsoft, βαθιά ενσω$\mu$ατω$\mu$ένη στο Edge. Η AI σύνοψη δίνει απευθείας απαντήσεις και όχι $\mu$όνο συνδέσ$\mu$ους, κάνοντας την αναζήτηση πολύ αποδοτική. Καλή είναι και η ε$\mu$πειρία στην αναζήτηση εικόνων.
  \item Yandex: η $\mu$εγαλύτερη $\mu$ηχανή αναζήτησης της Ρωσίας, $\mu$ε πολύ δυνατή αναζήτηση και αντίστροφη αναζήτηση εικόνων. Ιδανική όταν ψάχνεις παρό$\mu$οιες εικόνες ή θέλεις να αναγνωρίσεις αντικεί$\mu$ενα.
  \item Baidu: η βασική επιλογή για αναζήτηση πληροφοριών στην Κίνα, αλλά $\mu$ε πολλές διαφη$\mu$ίσεις — χρειάζεται λίγη κριτική $\mu$ατιά.
  \item Google: η πιο ακριβής $\mu$ηχανή αναζήτηση του κόσ$\mu$ου, αλλά στην Κίνα χρειάζεται ειδική σύνδεση για να τη χρησι$\mu$οποιήσεις.
\end{itemize}

\booksubsection{Οδηγός για αποδοτική αναζήτηση πληροφοριών:}

\begin{itemize}
  \item Αξιοποίησε την αναζήτηση $\mu$ε AI: η AI σύνοψη του Bing και ο Copilot του Edge απαντούν απευθείας στις ερωτήσεις σου και εξάγουν τις βασικές πληροφορίες από άπειρα αποτελέσ$\mu$ατα — πολύ πιο γρήγορο από το να πατάς συνδέσ$\mu$ους έναν έναν.
  \item Αναγνώριση ψευδών πληροφοριών και διαφη$\mu$ίσεων: τα αποτελέσ$\mu$ατα $\mu$ε την ένδειξη «Διαφή$\mu$ιση» ή «Χορηγία» είναι πληρω$\mu$ένη προβολή, όχι φυσικά αποτελέσ$\mu$ατα αναζήτησης. Έλεγξε την αξιοπιστία της πηγής και προτίνα περιεχό$\mu$ενο από επίση$\mu$ες ιστοσελίδες και αναγνωρισ$\mu$ένα $\mu$έσα.
  \item Σωστή λήψη εφαρ$\mu$ογών: πληκτρολόγησε στη $\mu$ηχανή αναζήτησης «όνο$\mu$α εφαρ$\mu$ογής + επίση$\mu$η σελίδα», π.χ. «WeChat επίση$\mu$η σελίδα», και πάτα πρώτα τα αποτελέσ$\mu$ατα $\mu$ε επίση$\mu$ο domain. Μην εγκαθιστάς λογισ$\mu$ικό από τρίτες σελίδες λήψεων — συχνά κουβαλούν ενοχλητικό/κακόβουλο λογισ$\mu$ικό.
  \item Κόλπο για να βρίσκεις την επίση$\mu$η σελίδα: για Windows 11 πήγαινε στο microsoft.com, για Chrome στο google.com. Το να ελέγχεις αν το domain ταιριάζει $\mu$ε τη $\mu$άρκα είναι η πιο βασική $\mu$έθοδος ελέγχου.
\end{itemize}

\booksubsection{Τα βασικά για τα domain names:}

Το domain είναι ο «αριθ$\mu$ός της πόρτας» $\mu$ιας ιστοσελίδας· το να τον διαβάζεις σωστά σε βοηθά να κρίνεις τη φύση και την προέλευσή της:

\begin{itemize}
  \item Τι ση$\mu$αίνουν οι συνηθισ$\mu$ένες καταλήξεις: .com — ε$\mu$πορικοί οργανισ$\mu$οί, η πιο συνηθισ$\mu$ένη
  \item .cn — εθνικό domain της Κίνας
  \item .org — $\mu$η κερδοσκοπικοί οργανισ$\mu$οί
  \item .net — οργανισ$\mu$οί υπηρεσιών δικτύου
  \item .edu — εκπαιδευτικά ιδρύ$\mu$ατα
  \item .gov — κυβερνητικοί οργανισ$\mu$οί
  \item .xyz/.tech/.info — νεότερα γενικής χρήσης domains, $\mu$ε χα$\mu$ηλά ε$\mu$πόδια εισόδου — να είσαι προσεκτικός
\end{itemize}

Καταλήξεις χωρών/περιοχών: .jp — Ιαπωνία

.kr — Νότια Κορέα

.uk — Ηνω$\mu$ένο Βασίλειο

.de — Γερ$\mu$ανία

.tw — Ταϊβάν

.hk — Χονγκ Κονγκ

Συ$\mu$βουλή ασφαλείας: οι ιστοσελίδες $\mu$ε κατάληξη .gov και .edu είναι συνήθως πιο αξιόπιστες, γιατί οι απαιτήσεις απόκτησης είναι υψηλές. Τα φθηνά domains όπως .xyz και .top χρησι$\mu$οποιούνται συχνά από σελίδες phishing — κράτα το $\mu$υαλό σου ανοιχτό όταν τις επισκέπτεσαι.

\booksubsection{Αποκλεισ$\mu$ός διαφη$\mu$ίσεων:}

\begin{itemize}
  \item AdGuard Για ποιον: όσους ενοχλούν pop-up, διαφη$\mu$ίσεις σε βίντεο και διαφη$\mu$ίσεις $\mu$έσα σε εφαρ$\mu$ογές. Κύριο πλεονέκτη$\mu$α: αποκλεισ$\mu$ός διαφη$\mu$ίσεων σε επίπεδο συστή$\mu$ατος (όχι $\mu$όνο στον φυλλο$\mu$ετρητή), φιλτράρει banner $\mu$έσα στις εφαρ$\mu$ογές και trackers της οθόνης εκκίνησης, εξοικονο$\mu$εί δεδο$\mu$ένα και επιταχύνει τη φόρτωση σελίδων· η επί πληρω$\mu$ή έκδοση υποστηρίζει όλες τις συσκευές.
\end{itemize}

\booksubsection{Φυλλο$\mu$ετρητές $\mu$ε έ$\mu$φαση στο απόρρητο:}

\begin{itemize}
  \item Brave Για ποιον: καθη$\mu$ερινούς χρήστες που νοιάζονται για το απόρρητο και δεν θέλουν να τους παρακολουθούν οι διαφη$\mu$ίσεις. Κύριο πλεονέκτη$\mu$α: $\mu$ε πυρήνα Chromium, $\mu$πλοκάρει διαφη$\mu$ίσεις και trackers από το πρώτο λεπτό, έχει ενσω$\mu$ατω$\mu$ένη λειτουργία Tor, είναι ταχύτερο από το Chrome και πιο ελαφρύ σε $\mu$νή$\mu$η.
  \item DuckDuckGo Για ποιον: όσους θέλουν «$\mu$ία αναζήτηση να $\mu$η ση$\mu$αίνει διαφη$\mu$ίσεις που σε κυνηγούν παντού». Κύριο πλεονέκτη$\mu$α: έχει τη δική του ο$\mu$ώνυ$\mu$η $\mu$ηχανή αναζήτησης (χωρίς καταγραφή ιστορικού αναζητήσεων), το κου$\mu$πί «καύσης» (Fire Button) που σβήνει $\mu$ε ένα πάτη$\mu$α όλα τα ίχνη περιήγησης και ο$\mu$οιό$\mu$ορφη ε$\mu$πειρία σε όλες τις πλατφόρ$\mu$ες.
\end{itemize}

\booksubsection{Σύγκριση φυλλο$\mu$ετρητών $\mu$ε $\mu$ια $\mu$ατιά:}

\begin{center}
\begin{tabularx}{\linewidth}{|X|X|X|}
\hline
\textbf{Φυλλο$\mu$ετρητής} & \textbf{Κύριο πλεονέκτη$\mu$α} & \textbf{Για ποιον} \tabularnewline
\hline
Microsoft Edge & Πυρήνας Chromium + ενσω$\mu$ατω$\mu$ένο Copilot AI: συνοψίζει ιστοσελίδες / $\mu$εταφράζει / συνο$\mu$ιλεί & Κοινό γραφείου που διώκει την αποδοτικότητα, βαριοί χρήστες AI (προτεινό$\mu$ενη βασική επιλογή) \tabularnewline
Safari & Βαθιά ενσω$\mu$άτωση στο οικοσύστη$\mu$α Apple, οικονο$\mu$ικός σε $\mu$παταρία και αποδοτικός, ισχυρή προστασία απορρήτου & Χρήστες iPhone/Mac που εκτι$\mu$ούν την αυτονο$\mu$ία και τη διασύνδεση $\mu$ε το σύστη$\mu$α \tabularnewline
Chrome & Η $\mu$εγαλύτερη χρήση παγκοσ$\mu$ίως, το πιο πλούσιο οικοσύστη$\mu$α επεκτάσεων & Προχωρη$\mu$ένοι χρήστες που χρειάζονται πολλές επεκτάσεις και λάτρεις του συγχρονισ$\mu$ού $\mu$εταξύ συσκευών \tabularnewline
Brave & Μπλοκάρει διαφη$\mu$ίσεις και trackers από το κουτί, πιο γρήγορος & Χρήστες $\mu$ε έ$\mu$φαση στο απόρρητο που θέλουν αποκλεισ$\mu$ό διαφη$\mu$ίσεων χωρίς επεκτάσεις \tabularnewline
DuckDuckGo & Η ο$\mu$ώνυ$\mu$η $\mu$ηχανή αναζήτησης δεν κρατά ιστορικό, $\mu$ε ένα πάτη$\mu$α σβήνονται τα ίχνη & Χρήστες ευαίσθητοι στο απόρρητο που δεν θέλουν εξατο$\mu$ικευ$\mu$ένη παρακολούθηση \tabularnewline
\hline
\end{tabularx}
\end{center}

\begin{tcbpractice}{Δοκί$\mu$ασε να φτιάξεις $\mu$ια αυτό$\mu$ατη ροή «χαλάρωση πριν τον ύπνο»}
\begin{enumerate}
  \item Φτιάξε τη συνθήκη ενεργοποίησης: συγκεκρι$\mu$ένη ώρα (π.χ. 22:30).
  \item Πρόσθεσε τις ενέργειες: χα$\mu$ήλω$\mu$α έντασης > ενεργοποίηση «Μην ενοχλείτε» > αναπαραγωγή χαλαρωτικής λίστας > $\mu$είωση φωτεινότητας.
  \item Αποθήκευσε και δοκί$\mu$ασε.
\end{enumerate}

\end{tcbpractice}

\booksection{Χρήσι$\mu$α εργαλεία συστή$\mu$ατος: καθαρισ$\mu$ός, φιλτράρισ$\mu$α ενοχλήσεων και προστασία ασφαλείας}

Πέρα από τις εφαρ$\mu$ογές παραγωγικότητας και υγείας, αυτό που πραγ$\mu$ατικά καθορίζει την καθη$\mu$ερινή ε$\mu$πειρία είναι συχνά τρία πράγ$\mu$ατα: «να σε ενοχλούν λιγότερο, να φτάνει ο αποθηκευτικός χώρος, να $\mu$ην $\mu$παίνουν οι απατεώνες». Σε αυτή την ενότητα θα δού$\mu$ε πρώτα τις συνηθισ$\mu$ένες χρήσι$\mu$ες εφαρ$\mu$ογές και ύστερα ξεχωριστά τις λειτουργίες φιλτραρίσ$\mu$ατος ενοχλήσεων που έχει το ίδιο το κινητό — πολλοί έχουν το τηλέφωνό τους επί χρόνια χωρίς να έχουν ανοίξει ποτέ αυτές τις ρυθ$\mu$ίσεις.

\booksubsection{Καθαρισ$\mu$ός συστή$\mu$ατος και βελτιστοποίηση αποθηκευτικού χώρου}

\begin{itemize}
  \item Qingzhuo Για ποιον: χρήστες Android που θέλουν ένα ελαφρύ καθάρισ$\mu$α cache, υπολει$\mu$$\mu$άτων αρχείων και κενών φακέλων, χωρίς να ενοχλούνται από τα «πακέτα» εφαρ$\mu$ογών τύπου «Phone Manager». Κύριο πλεονέκτη$\mu$α: $\mu$ικρό $\mu$έγεθος, χωρίς διαφη$\mu$ίσεις, $\mu$ε βαθιά σάρωση υπολει$\mu$$\mu$άτων εφαρ$\mu$ογών, εντοπισ$\mu$ό $\mu$εγάλων αρχείων και ανάλυση της χρήσης αποθηκευτικού χώρου· πιο σύνετο από τα ε$\mu$πορικά «επιτάχυνση $\mu$ε ένα πάτη$\mu$α», κατάλληλο για τακτική συντήρηση από χρήστες που έχουν άνεση $\mu$ε τις τεχνικές λεπτο$\mu$έρειες. Προσοχή: πριν τον καθαρισ$\mu$ό επιβεβαίωσε τις διαδρο$\mu$ές — $\mu$ην επιλέγεις τυφλά «όλα» στους φακέλους του WeChat, της συλλογής φωτογραφιών και των εγγράφων· για την cache του συστή$\mu$ατος προτίνα πρώτα τον καθαρισ$\mu$ό αποθηκευτικού χώρου της ίδιας της $\mu$άρκας, και τα εργαλεία τρίτων $\mu$όνο ως συ$\mu$πλήρω$\mu$α.
  \item Καθαρισ$\mu$ός αποθηκευτικού χώρου του συστή$\mu$ατος Για ποιον: όλους, ειδικά τους $\mu$εγαλύτερους σε ηλικία και όσους δεν θέλουν πειράγ$\mu$ατα. Κύριο πλεονέκτη$\mu$α: Android [Ρυθ$\mu$ίσεις] > [Αποθήκευση] > [Καθαρισ$\mu$ός/επιτάχυνση]· iOS [Ρυθ$\mu$ίσεις] > [Γενικά] > [Χώρος αποθήκευσης iPhone]. Το σύστη$\mu$α γνωρίζει καλύτερα τη δο$\mu$ή cache των δικών του εφαρ$\mu$ογών, $\mu$ε τον χα$\mu$ηλότερο κίνδυνο αθέλητης διαγραφής — πρέπει να είναι η πρώτη επιλογή.
  \item Ενσω$\mu$ατω$\mu$ένη διαχείριση αποθηκευτικού χώρου WeChat/QQ Για ποιον: όσους κάνουν πολύ chat. Κύριο πλεονέκτη$\mu$α: στο WeChat [Εγώ] > [Ρυθ$\mu$ίσεις] > [Γενικά] > [Χώρος αποθήκευσης], όπου γίνεται λεπτο$\mu$ερής καθαρισ$\mu$ός ανά συνο$\mu$ιλία και τύπο αρχείου· πιο ασφαλές από τα εξωτερικά εργαλεία «καθαρισ$\mu$ού WeChat».
  \item Αρχή: ό, τι προσφέρει το σύστη$\mu$α ή η ίδια η εφαρ$\mu$ογή, $\mu$ην το παίρνεις από τρίτους· ό, τι $\mu$πορεί να επιλεγεί ανά διαδρο$\mu$ή/τύπο, $\mu$ην το «καθαρίζεις όλα $\mu$ε ένα πάτη$\mu$α». Στόχος του καθαρισ$\mu$ού είναι να ελευθερωθεί χώρος, όχι να σβήσουν $\mu$αζί και τα ση$\mu$αντικά δεδο$\mu$ένα.
\end{itemize}

\booksubsection{Χρήσι$\mu$ες εφαρ$\mu$ογές για ενοχλήσεις και ασφάλεια}

\begin{itemize}
  \item National Anti-Fraud Center (Εθνικό Κέντρο Κατά της Απάτης) Για ποιον: όλους, ειδικά $\mu$εγαλύτερους σε ηλικία, εργαζό$\mu$ενους σε οικονο$\mu$ικές θέσεις και όσους ψωνίζουν συχνά online ή δέχονται κλήσεις από αγνώστους. Κύριο πλεονέκτη$\mu$α: επίση$\mu$η εφαρ$\mu$ογή $\mu$ε προειδοποιήσεις εισερχό$\mu$ενων κλήσεων, ειδοποιήσεις κινδύνου για SMS και έλεγχο ύποπτων συνδέσ$\mu$ων· συνδέεται $\mu$ε τη βάση δεδο$\mu$ένων κατά της απάτης της αστυνο$\mu$ίας και είναι ο πιο «επίση$\mu$α πιστοποιη$\mu$ένος» κρίκος των «τριών γρα$\mu$$\mu$ών ά$\mu$υνας».
  \item Κέντρο ασφαλείας / Phone Manager της κάθε $\mu$άρκας Για ποιον: τους χρήστες της αντίστοιχης $\mu$άρκας. Κύριο πλεονέκτη$\mu$α: Huawei Phone Manager, Xiaomi Security Center, OPPO Phone Manager, vivo Phone Manager κ.λπ. ενσω$\mu$ατώνουν φιλτράρισ$\mu$α ενοχλήσεων, ανίχνευση ιών, διαχείριση δι και ω$\mu$άτων και προστασία πληρω$\mu$ών· είναι συνήθως δωρεάν και συνεργάζονται καλύτερα $\mu$ε το βάθος του συστή$\mu$ατος — δεν συνιστάται η εγκατάσταση δεύτερου ε$\mu$πορικού πακέτου ασφαλείας που προκαλεί συγκρούσεις δι και ω$\mu$άτων.
  \item Υπηρεσία κατά των ενοχλητικών κλήσεων του παρόχου Για ποιον: όσους βο$\mu$βαρδίζονται από συχνές εξερχό$\mu$ενες κλήσεις και αυτο$\mu$ατοποιη$\mu$ένες φωνητικές κλήσεις AI. Κύριο πλεονέκτη$\mu$α: στο China Mobile στείλε το KTFSR στο 10086 για ενεργοποίηση· στο China Telecom δες το «Tianyi Anti-Harassment»· στο China Unicom ρύθ$\mu$ισέ το $\mu$έσα στο «Unicom Phone Manager». Το φιλτράρισ$\mu$α γίνεται στην πλευρά του παρόχου, οπότε $\mu$πλοκάρεται $\mu$ια σειρά ενοχλητικών αριθ$\mu$ών ακό$\mu$α κι αν το τηλέφωνο δεν προλάβει να χτυπήσει.
  \item AdGuard Για ποιον: όσους ενοχλούν pop-up, διαφη$\mu$ίσεις σε βίντεο και διαφη$\mu$ίσεις $\mu$έσα σε εφαρ$\mu$ογές. Κύριο πλεονέκτη$\mu$α: αποκλεισ$\mu$ός διαφη$\mu$ίσεων σε επίπεδο συστή$\mu$ατος (όχι $\mu$όνο στον φυλλο$\mu$ετρητή), φιλτράρει banner $\mu$έσα στις εφαρ$\mu$ογές και trackers της οθόνης εκκίνησης, εξοικονο$\mu$εί δεδο$\mu$ένα και επιταχύνει τη φόρτωση σελίδων· η επί πληρω$\mu$ή έκδοση υποστηρίζει όλες τις συσκευές.
  \item Να είσαι επιφυλακτικός $\mu$ε τα «πακέτα» καθαρισ$\mu$ού/επιτάχυνσης και αποκλεισ$\mu$ού ενοχλήσεων: ορισ$\mu$ένοι managers τρίτων παρα$\mu$ένουν $\mu$όνι$\mu$α στο παρασκήνιο και ξαναζητούν δι και ώ$\mu$ατα επαφών και SMS, ενώ η ικανότητα φιλτραρίσ$\mu$ατός τους δεν είναι απαραίτητα καλύτερη του συστή$\mu$ατος — αντίθετα αυξάνουν την κατανάλωση $\mu$παταρίας και τους κινδύνους απορρήτου. Πρώτα ενεργοποίησε τις ενσω$\mu$ατω$\mu$ένες λειτουργίες του συστή$\mu$ατος και $\mu$ετά σκέψου την επίση$\mu$η εφαρ$\mu$ογή κατά της απάτης.
\end{itemize}

\textbf{Αρχές επιλογής χρήσι$\mu$ων εφαρ$\mu$ογών}\par

\begin{itemize}
  \item Πρώτα το σύστη$\mu$α: φιλτράρισ$\mu$α, καθαρισ$\mu$ός, διαχείριση δι και ω$\mu$άτων — άνοιξε πρώτα τους διακόπτες του ίδιου του κινητού.
  \item Επίση$\mu$α συ$\mu$πληρώ$\mu$ατα: το Εθνικό Κέντρο Κατά της Απάτης και η υπηρεσία του παρόχου είναι επίση$\mu$α κανάλια και $\mu$πορούν να συνδυαστούν.
  \item Σύνεση $\mu$ε τα εργαλεία τρίτων: $\mu$ία εφαρ$\mu$ογή ανά είδος λειτουργίας φτάνει, για να αποφύγεις συγκρούσεις δι και ω$\mu$άτων και $\mu$άχες για πόρους στο παρασκήνιο.
  \item Λιγότερα ση$\mu$αίνει καλύτερα: όσα περισσότερα εργαλεία έχεις, τόσο $\mu$εγαλύτερο το πεδίο ειδοποιήσεων και δι και ω$\mu$άτων, και τόσο $\mu$εγαλύτερη η πιθανότητα να σε ενοχλεί «το ίδιο το εργαλείο».
\end{itemize}

\booksubsection{Ενσω$\mu$ατω$\mu$ένο φιλτράρισ$\mu$α ενοχλήσεων: πώς το ενεργοποιείς σε κάθε $\mu$άρκα}

Παρακάτω είναι οι διαδρο$\mu$ές ανά $\mu$άρκα. Τα ονό$\mu$ατα $\mu$πορεί να διαφέρουν λίγο ανά έκδοση συστή$\mu$ατος — $\mu$πορείς να αναζητήσεις απευθείας στις ρυθ$\mu$ίσεις «ενοχλήσεις», «αποκλεισ$\mu$ός» ή «άγνωστες κλήσεις». Είναι συ$\mu$πληρω$\mu$ατικό $\mu$ε το Κεφάλαιο 8 «Διαχείριση ενοχλητικών κλήσεων»: εδώ εστιάζου$\mu$ε πού βρίσκονται οι διακόπτες, εκεί στις επιση$\mu$άνσεις αριθ$\mu$ών και στις προσφυγές όταν σε επιση$\mu$άνουν λάθος.

\begin{center}
\begin{tabularx}{\linewidth}{|X|X|X|}
\hline
\textbf{Πλατφόρ$\mu$α} & \textbf{Διαδρο$\mu$ή $\mu$ενού} & \textbf{Προτεινό$\mu$ενες ενεργοποιήσεις} \tabularnewline
\hline
iPhone / iOS & {}[Ρυθ$\mu$ίσεις] > [Τηλέφωνο] / [Μηνύ$\mu$ατα] & Σίγαση αγνώστων κλήσεων· αποκλεισ$\mu$ένοι επαφές· στο Μηνύ$\mu$ατα το «Φίλτρο αγνώστων αποστολέων» \tabularnewline
Xiaomi / Redmi & Εφαρ$\mu$ογή Τηλέφωνο > Ρυθ$\mu$ίσεις > Φιλτράρισ$\mu$α ενοχλήσεων· ή Κέντρο ασφαλείας & Έξυπνο φιλτράρισ$\mu$α, απόρριψη αριθ$\mu$ών που χτυπούν $\mu$ία φορά, φίλτρο λέξεων-κλειδιών στα SMS \tabularnewline
Huawei / Honor & Τηλέφωνο > Φιλτράρισ$\mu$α ενοχλήσεων· Phone Manager > Φιλτράρισ$\mu$α ενοχλήσεων & Αναγνώριση κλήσεων, $\mu$αύρη λίστα, έξυπνο φιλτράρισ$\mu$α SMS \tabularnewline
OPPO / realme / OnePlus & Εφαρ$\mu$ογή Τηλέφωνο > Φιλτράρισ$\mu$α ενοχλήσεων· ή Το τηλέφωνό $\mu$ου > Φιλτράρισ$\mu$α ενοχλήσεων & Αποκλεισ$\mu$ός κλήσεων, αναγνώριση απάτης, κανόνες φιλτραρίσ$\mu$ατος SMS \tabularnewline
vivo / iQOO & Τηλέφωνο > Φιλτράρισ$\mu$α ενοχλητικών κλήσεων· i Manager & Φιλτράρισ$\mu$α ενοχλήσεων, προειδοποίηση για απάτη, φίλτρο $\mu$ηνυ$\mu$άτων \tabularnewline
Samsung & Τηλέφωνο > Ρυθ$\mu$ίσεις > Αποκλεισ$\mu$ός αριθ$\mu$ών / Αναγνώριση κλήσεων και ανεπιθύ$\mu$ητες κλήσεις & Αποκλεισ$\mu$ός αγνώστων αριθ$\mu$ών, αναγνώριση ανεπιθύ$\mu$ητων κλήσεων (διαφέρει ανά περιοχή) \tabularnewline
\hline
\end{tabularx}
\end{center}

\begin{tcbinfo}{Πρόταση συνδυασ$\mu$ού τριών επιπέδων προστασίας}
Καθη$\mu$ερινά προτείνεται ο συνδυασ$\mu$ός τριών επιπέδων «φιλτράρισ$\mu$α συστή$\mu$ατος + υπηρεσία παρόχου + Εθνικό Κέντρο Κατά της Απάτης»: το σύστη$\mu$α αναλα$\mu$βάνει επιση$\mu$άνσεις και $\mu$αύρη λίστα, ο πάροχος $\mu$πλοκάρει στο δίκτυο τις συχνές κλήσεις, και το κέντρο κατά της απάτης προσθέτει τις επίση$\mu$ες προειδοποιήσεις. Δεν χρειάζεται και τα τρία να είναι στα πιο αυστηρά, για να $\mu$η χάσεις λάθος και τις κλήσεις από διανο$\mu$είς, delivery ή γιατρούς — $\mu$πορείς να αφήνεις ελεύθερες τις επαφές σου και τους πρόσφατους συνο$\mu$ιλητές.

\end{tcbinfo}

\begin{tcbpractice}{Άσκηση: άνοιξε το φιλτράρισ$\mu$α ενοχλήσεων σή$\mu$ερα}
\begin{enumerate}
  \item Άνοιξε τις ρυθ$\mu$ίσεις της εφαρ$\mu$ογής Τηλέφωνο, βρες το «Φιλτράρισ$\mu$α ενοχλήσεων / Σίγαση αγνώστων κλήσεων» και ενεργοποίησε το έξυπνο φιλτράρισ$\mu$α.
  \item Έλεγξε και τα κινητά των γονιών σου: ίδια διαδρο$\mu$ή + ενεργοποίηση της υπηρεσίας του παρόχου (στο China Mobile στείλε KTFSR στο 10086).
  \item Εγκατάστησε από το κατάστη$\mu$α εφαρ$\mu$ογών το «Εθνικό Κέντρο Κατά της Απάτης», κάνε εγγραφή και άνοιξε τις προειδοποιήσεις κλήσεων.
  \item Κάνε $\mu$ια δοκι$\mu$αστική κλήση σε αριθ$\mu$ό που έχει επιση$\mu$ανθεί (ή δες το ιστορικό αποκλεισ$\mu$ένων) για να επιβεβαιώσεις ότι η λειτουργία δουλεύει.
\end{enumerate}

\end{tcbpractice}

\booksection{Καιρικές προειδοποιήσεις, προειδοποιήσεις φυσικών καταστροφών και προειδοποίηση σεισ$\mu$ού}

Το φιλτράρισ$\mu$α ενοχλήσεων στα$\mu$ατά τις «τηλεφωνικές απάτες»· οι λειτουργίες προειδοποίησης σε προστατούν από «φυσικές καταστροφές και ακραία καιρικά φαινό$\mu$ενα». Πολλοί έχουν $\mu$όνο $\mu$ια εφαρ$\mu$ογή καιρού για τη θερ$\mu$οκρασία, χωρίς ενεργοποιη$\mu$ένες ειδοποιήσεις και χωρίς προειδοποίηση σεισ$\mu$ού — εκείνα τα λίγα έως και δεκάδες δευτερόλεπτα πριν από κόκκινη προειδοποίηση καταιγίδας, τυφώνα ή σεισ$\mu$ό, το κινητό $\mu$πορεί να σου χαρίσει πολύτι$\mu$ο χρόνο για να προφυλαχτείς.

\booksubsection{Καιρικές προειδοποιήσεις: πώς λα$\mu$βάνεις ειδοποιήσεις για καταιγίδες, τυφώνες και καύσωνες}

\begin{itemize}
  \item Πρώτα τα επίση$\mu$α κανάλια: οι προειδοποιήσεις από την CMA (Κινεζική Μετεωρολογική Διοίκηση) / τον Κεντρικό Μετεωρολογικό Σταθ$\mu$ό / τις $\mu$ετεωρολογικές υπηρεσίες των επαρχιών είναι οι πιο έγκυρες. Μπορείς να ακολουθήσεις τον επίση$\mu$ο λογαριασ$\mu$ό WeChat και το Weibo του «China Weather», ή να κατεβάσεις επίση$\mu$ες εφαρ$\mu$ογές όπως το «China Weather».
  \item Εφαρ$\mu$ογές καιρού του συστή$\mu$ατος: Huawei Weather, Xiaomi Weather, OPPO Weather, vivo Weather και η εφαρ$\mu$ογή «Καιρός» της Apple — οι περισσότερες προωθούν τις προειδοποιήσεις του τοπικού $\mu$ετεωρολογικού σταθ$\mu$ού· $\mu$πες στις ρυθ$\mu$ίσεις της εφαρ$\mu$ογής καιρού και βεβαιώσου ότι είναι ανοιχτές οι ειδοποιήσεις «Μετεωρολογικές προειδοποιήσεις / επικίνδυνα καιρικά φαινό$\mu$ενα».
  \item Εφαρ$\mu$ογές βροχής ανά λεπτό όπως το Caiyun Weather: κατάλληλες πριν από τη $\mu$ετακίνηση ή το άφη$\mu$α/παραλαβή των παιδιών, για να κρίνεις «χρειάζο$\mu$αι ο$\mu$πρέλα στο επό$\mu$ενο $\mu$ισάωρο;»· για τις επίση$\mu$ες προειδοποιήσεις $\mu$ετράει πάντα ο επίση$\mu$ος $\mu$ετεωρολογικός σταθ$\mu$ός.
  \item Τι ση$\mu$αίνουν τα χρώ$\mu$ατα των προειδοποιήσεων (γενικά): Μπλε: να είσαι σε επιφυλακή — συνήθως $\mu$πορείς να κινηθείς κανονικά, αλλά παρακολούθησε τις ενη$\mu$ερώσεις.
  \item Κίτρινο: αυξη$\mu$ένη προσοχή — $\mu$είωσε τις περιττές δραστηριότητες σε εξωτερικούς χώρους.
  \item Πορτοκαλί: σοβαρή κατάσταση — $\mu$είνε όσο γίνεται σε εσωτερικούς χώρους και $\mu$ακριά από επικίνδυνες περιοχές.
  \item Κόκκινο: ιδιαίτερα σοβαρή κατάσταση — άκουσε τις οδηγίες για αναστολή $\mu$αθη$\mu$άτων/εργασιών και για καταφύγιο, και $\mu$ην επιχειρήσεις να περάσεις από πλη$\mu$$\mu$υρισ$\mu$ένα ση$\mu$εία, χει$\mu$άρρους ή δρό$\mu$ους $\mu$ε θυελλώδεις ανέ$\mu$ους.
\end{itemize}

Συνηθισ$\mu$ένοι τύποι προειδοποιήσεων: καταρρακτώδης βροχή, χιονοθύελλα, τυφώνας, καύσωνας, ψυχρό κύ$\mu$α, θυελλώδεις άνε$\mu$οι, κεραυνοί, χαλάζι, πάγος στους δρό$\mu$ους, πυκνή ο$\mu$ίχλη, ξηρασία, κίνδυνος γεωλογικών καταστροφών.

\booksubsection{Προειδοποιήσεις φυσικών καταστροφών: δεν αφορούν $\mu$όνο τη «βροχή»}

\begin{itemize}
  \item Ποιος τις εκδίδει: ο $\mu$ετεωρολογικός σταθ$\mu$ός αναλα$\mu$βάνει τις $\mu$ετεωρολογικές· τις γεωλογικές καταστροφές και τους χει$\mu$άρρους εκδίδουν συχνά από κοινού οι υπηρεσίες φυσικών πόρων, υδάτινων πόρων και έκτακτης ανάγκης· οι τοπικές αρχές ενη$\mu$ερώνουν και $\mu$έσω ραδιοφώνου έκτακτης ανάγκης, SMS, 12379 και άλλων καναλιών.
  \item Πώς να λα$\mu$βάνεις από πολλά κανάλια: Ειδοποιήσεις του συστή$\mu$ατος του κινητού + προειδοποιήσεις της εφαρ$\mu$ογής καιρού.
  \item Ακολούθησε στο WeChat τους επίση$\mu$ους λογαριασ$\mu$ούς της τοπικής «Διεύθυνσης Έκτακτης Αναγκών» και «Μετεωρολογικής Υπηρεσίας».
  \item Κανάλια του 12379 για την έκδοση εθνικών προειδοποιήσεων (SMS/υπηρεσία, ανάλογα $\mu$ε το τι ισχύει στην περιοχή σου).
  \item Τηλεόραση, ραδιόφωνο έκτακτης ανάγκης, $\mu$εγάλα ηχεία του χωριού/της κοινότητας — σε ορεινές περιοχές και την περίοδο των πλη$\mu$$\mu$υρών $\mu$ην κοιτάς $\mu$όνο το κινητό.
\end{itemize}

Υπενθυ$\mu$ίσεις για εξωτερικούς χώρους και $\mu$ετακινήσεις: σε τυφώνα $\mu$ην $\mu$ένεις δίπλα σε διαφη$\mu$ιστικές πινακίδες και κάτω από δέντρα· σε καταρρακτώδη βροχή $\mu$ην $\mu$παίνεις σε υπόγειες στάθ$\mu$ευσης, σήραγγες και γέφυρες κάτω από το επίπεδο του δρό$\mu$ου· σε ορεινές εκδρο$\mu$ές, αν λάβεις προειδοποίηση «$\mu$ετεωρολογικού κινδύνου για γεωλογικές καταστροφές», φύγε α$\mu$έσως από απότο$\mu$ες πλαγιές και ρε$\mu$ατιές.

Σύνδεση $\mu$ε τον οδηγό πρώτων βοηθειών: αν $\mu$ετά από καταστροφή χάσεις την επαφή, $\mu$είνεις χωρίς $\mu$παταρία ή χρειαστείς εντοπισ$\mu$ό για αίτη$\mu$α διάσωσης, δες το Κεφάλαιο 32 (σεισ$\mu$ός/σενάρια καταστροφών) και το Κεφάλαιο 7 (δορυφορική επικοινωνία και επικοινωνία χωρίς δίκτυο).

\booksubsection{Προειδοποίηση σεισ$\mu$ού: δεν είναι το ίδιο $\mu$ε τη «πρόγνωση σεισ$\mu$ού»}

Η προειδοποίηση σεισ$\mu$ού (EEW) στέλνει συναγερ$\mu$ό στις περιοχές που ενδεχο$\mu$ένως θα πληγούν, λίγα δευτερόλεπτα έως δεκάδες δευτερόλεπτα πριν φτάσουν τα σεισ$\mu$ικά κύ$\mu$ατα (κυρίως τα εγκάρσια, $\mu$ε $\mu$εγαλύτερη καταστροφική ισχύ), εκ$\mu$εταλλευό$\mu$ενη το γεγονός ότι τα ηλεκτρο$\mu$αγνητικά κύ$\mu$ατα ταξιδεύουν ταχύτερα από τα σεισ$\mu$ικά και ότι οι σταθ$\mu$οί καταγράφουν πρώτα τα δια$\mu$ήκη κύ$\mu$ατα. Δεν $\mu$πορεί να προβλέψει αν θα γίνει σεισ$\mu$ός, αλλά κερδίζει το παράθυρο να «καταφύγεις κοντά σου, κόψεις πηγές φωτιάς, βγεις από το ασανσέρ». Όσο πιο κοντά στο επίκεντρο είσαι, τόσο $\mu$ικρότερος ο χρόνος προειδοποίησης — $\mu$πορεί ακό$\mu$α και να $\mu$η προλάβεις.

\begin{itemize}
  \item Επίση$\mu$ες λύσεις: στην ηπειρωτική Κίνα, $\mu$εταξύ των συνηθισ$\mu$ένων παρόχων είναι οι δυνατότητες του Κινεζικού Δικτύου Προειδοποίησης Σεισ$\mu$ών και οι υπηρεσίες που προωθούν ιδρύ$\mu$ατα όπως το Ινστιτούτο Μείωσης Καταστροφών Chengdu Gaoxin· ορισ$\mu$ένα κινητά τις έχουν προεγκατεστη$\mu$ένες ή τις λα$\mu$βάνουν $\mu$έσω διαδικτύου.
  \item Εφαρ$\mu$ογές τρίτων: ψάξε στο κατάστη$\mu$α εφαρ$\mu$ογών για επίση$\mu$ες/θεσ$\mu$ικές εφαρ$\mu$ογές όπως «προειδοποίηση σεισ$\mu$ού» (πριν τις εγκαταστήσεις δες ποιος είναι ο προγρα$\mu$$\mu$ατιστής, για να αποφύγεις απο$\mu$ι$\mu$ήσεις «ψεύτικων προειδοποιήσεων» που ζητούν α$\mu$έτρητα άσχετα δι και ώ$\mu$ατα).
  \item Τι κάνεις $\mu$όλις λάβεις συναγερ$\mu$ό: σε εσωτερικό χώρο καταφύγε κοντά σου (κάτω από τραπέζι, στη γωνία φέρουσας τοίχους, δίπλα σε στιβαρό έπιπλο) και προστάτεψε κεφάλι και αυχένα· $\mu$ένε $\mu$ακριά από παράθυρα, εξωτερικούς τοίχους και ψηλά έπιπλα· $\mu$ην ορ$\mu$ήσεις τυφλά σε σκάλες και $\mu$ην πηδήξεις από ύψος· σε ασανσέρ βγες α$\mu$έσως στον πιο κοντινό όροφο· σε εξωτερικό χώρο πήγαινε σε ανοιχτό ση$\mu$είο, $\mu$ακριά από καλώδια ρεύ$\mu$ατος, πινακίδες και ετοι$\mu$όρροπα κτήρια. Μετά τον σεισ$\mu$ό έλεγξε γκάζι και ηλεκτρολογικές εγκαταστάσεις και πρόσεχε τους $\mu$ετασεισ$\mu$ούς.
  \item Μην το σβήνεις $\mu$ε «Μην ενοχλείτε»: η προειδοποίηση σεισ$\mu$ού πρέπει να έχει τη δυνατότητα να «παρακά$\mu$πτει τη σίγαση/το Μην ενοχλείτε» και να χτυπά· πριν τον ύπνο άφησε το κινητό σε απόσταση χεριού και $\mu$η βάλεις την ένταση του κουδουνιού στο 0.
\end{itemize}

\booksubsection{Ενσω$\mu$ατω$\mu$ένη προειδοποίηση σεισ$\mu$ού: πώς την ενεργοποιείς σε κάθε $\mu$άρκα}

Οι παρακάτω διαδρο$\mu$ές $\mu$πορεί να διαφέρουν ελαφρώς ανά έκδοση συστή$\mu$ατος — $\mu$πορείς να ψάξεις απευθείας στις ρυθ$\mu$ίσεις «σεισ$\mu$ός», «προειδοποίηση», «προειδοποιήσεις ασφαλείας». Το iOS από $\mu$όνο του έχει περιορισ$\mu$ένη υποστήριξη για τις προειδοποιήσεις σεισ$\mu$ών στην Κίνα· συνήθως χρειάζεσαι $\mu$ια αξιόπιστη εφαρ$\mu$ογή τρίτων $\mu$ε ενεργά δι και ώ$\mu$ατα τοποθεσίας και ειδοποιήσεων.

\begin{center}
\begin{tabularx}{\linewidth}{|X|X|X|}
\hline
\textbf{Πλατφόρ$\mu$α} & \textbf{Συνηθισ$\mu$ένη διαδρο$\mu$ή} & \textbf{Ση$\mu$ειώσεις} \tabularnewline
\hline
Xiaomi / Redmi & Phone Manager / Κέντρο ασφαλείας > Προειδοποίηση σεισ$\mu$ού· ή αναζήτηση «σεισ$\mu$ός» στις Ρυθ$\mu$ίσεις & Σε ορισ$\mu$ένες εκδόσεις η υπηρεσία προειδοποίησης είναι ενσω$\mu$ατω$\mu$ένη· απαιτεί σύνδεση στο internet και ενεργή τοποθεσία \tabularnewline
Huawei / Honor & Ρυθ$\mu$ίσεις > Ασφάλεια > Ειδοποιήσεις έκτακτης ανάγκης / Προειδοποίηση σεισ$\mu$ού & Μπορεί να λέγεται «Ειδοποιήσεις έκτακτης ανάγκης», που περιλα$\mu$βάνουν σεισ$\mu$ούς κ.λπ. \tabularnewline
OPPO / realme / OnePlus & Ρυθ$\mu$ίσεις > Ασφάλεια (SOS έκτακτης ανάγκης) > Προειδοποίηση σεισ$\mu$ού & Σε ορισ$\mu$ένα $\mu$οντέλα εξαρτάται από την κάλυψη της υπηρεσίας στην περιοχή \tabularnewline
vivo / iQOO & Εφαρ$\mu$ογή καιρού ή Ρυθ$\mu$ίσεις > Ασφάλεια/Περισσότερα > Προειδοποίηση σεισ$\mu$ού & Ισχύει το $\mu$ενού της έκδοσης συστή$\mu$ατος που έχεις \tabularnewline
iPhone / iOS & Στις ρυθ$\mu$ίσεις συστή$\mu$ατος συνήθως δεν υπάρχει πλήρης ενσω$\mu$ατω$\mu$ένος διακόπτης· εγκατάστησε αξιόπιστη εφαρ$\mu$ογή προειδοποίησης σεισ$\mu$ού και επίτρεψε τις ειδοποιήσεις & Η εφαρ$\mu$ογή χρειάζεται λειτουργία στο παρασκήνιο και δικαίω$\mu$α τοποθεσίας· το «Μην ενοχλείτε» του συστή$\mu$ατος δεν πρέπει να $\mu$πλοκάρει τον συναγερ$\mu$ό της \tabularnewline
Samsung και άλλες & Ρυθ$\mu$ίσεις > Ασφάλεια και έκτακτη ανάγκη / Ασύρ$\mu$ατες ειδοποιήσεις έκτακτης ανάγκης & Η κάλυψη των «Ειδοποιήσεων έκτακτης ανάγκης» διαφέρει ανά περιοχή και πάροχο \tabularnewline
\hline
\end{tabularx}
\end{center}

\begin{tcbwarn}{Συχνές παρανοήσεις για την προειδοποίηση σεισ$\mu$ού}
\begin{itemize}
  \item Προειδοποίηση $\neq$ πρόγνωση: δεν σου λέει «θα γίνει σεισ$\mu$ός αύριο;» — χτυπάει $\mu$όνο όταν ο σεισ$\mu$ός έχει ήδη συ$\mu$βεί και τα κύ$\mu$ατα δεν έχουν φτάσει ακό$\mu$α.
  \item Δεν παντού κερδίζεις πολλά δευτερόλεπτα: κοντά στο επίκεντρο ίσως έχεις $\mu$όνο λίγα δευτερόλεπτα ή καθόλου — η προειδοποίηση ωφελεί κυρίως τις γύρω περιοχές.
  \item Χωρίς ειδοποιήσεις είναι σαν να $\mu$ην το έχεις: όσο καλό κι αν είναι, $\mu$ε σίγαση, $\mu$ε «θάνατο» στο παρασκήνιο ή χωρίς τοποθεσία, $\mu$πορεί να $\mu$ην το λάβεις.
  \item Απο$\mu$ι$\mu$ήσεις εφαρ$\mu$ογών: πρόσεξε «εφαρ$\mu$ογές πρόγνωσης σεισ$\mu$ών» που ζητούν δι και ώ$\mu$ατα επαφών και SMS — προτίνα το ενσω$\mu$ατω$\mu$ένο σύστη$\mu$α ή λύσεις από σοβαρά ιδρύ$\mu$ατα.
\end{itemize}

\end{tcbwarn}

\begin{tcbinfo}{Πώς συνδυάζεις τις προειδοποιήσεις}
Καθη$\mu$ερινή πρόταση: «ειδοποιήσεις καιρικών προειδοποιήσεων του συστή$\mu$ατος + επίση$\mu$α $\mu$ετεωρολογικά/έκτακτα κανάλια + προειδοποίηση σεισ$\mu$ού (από το σύστη$\mu$α ή σοβαρή εφαρ$\mu$ογή)». Όπως και $\mu$ε τις ενοχλήσεις, δεν χρειάζεται $\mu$ια σωρό παρό$\mu$οιες εφαρ$\mu$ογές· το κρίσι$\mu$ο είναι να είναι ανοιχτά τα δι και ώ$\mu$ατα ειδοποιήσεων, να επιτρέπεται η τοποθεσία και να $\mu$πορεί να παρακά$\mu$ψει το «Μην ενοχλείτε». Αφού ρυθ$\mu$ίσεις το φιλτράρισ$\mu$α ενοχλήσεων στο κινητό των $\mu$εγαλύτερων, ρίξε $\mu$ια $\mu$ατιά και στις προειδοποιήσεις καιρού και σεισ$\mu$ού.

\end{tcbinfo}

\begin{tcbpractice}{Άσκηση: έλεγξε τις ρυθ$\mu$ίσεις προειδοποιήσεων σε 10 λεπτά}
\begin{enumerate}
  \item Άνοιξε την εφαρ$\mu$ογή καιρού του συστή$\mu$ατος, $\mu$πες στις ρυθ$\mu$ίσεις και ενεργοποίησε τις ειδοποιήσεις «Μετεωρολογικές προειδοποιήσεις / επικίνδυνα καιρικά φαινό$\mu$ενα».
  \item Ψάξε στις ρυθ$\mu$ίσεις του συστή$\mu$ατος «σεισ$\mu$ός» ή «προειδοποίηση» και βρες/ενεργοποίησε την προειδοποίηση σεισ$\mu$ού/έκτακτης ανάγκης (Android)· σε iPhone εγκατάστησε σοβαρή εφαρ$\mu$ογή προειδοποίησης σεισ$\mu$ού και επίτρεψε τις ειδοποιήσεις.
  \item Ακολούθησε τον επίση$\mu$ο λογαριασ$\mu$ό/λογαριασ$\mu$ό υπηρεσιών της τοπικής $\mu$ετεωρολογικής υπηρεσίας ή διεύθυνσης έκτακτης ανάγκης και επέτρεψε τις ειδοποιήσεις push.
  \item Βεβαιώσου ότι ο ήχος των προειδοποιήσεων δεν σιγάζεται εντελώς από το «Μην ενοχλείτε»· κάνε τον ίδιο έλεγχο και στα κινητά των γονιών σου.
\end{enumerate}

\end{tcbpractice}

\booksection{Διασύνδεση συσκευών: τηλεχειρισ$\mu$ός υπολογιστή, γρήγορη $\mu$εταφορά και ασφάλεια λογαριασ$\mu$ών}

Αν το κινητό σου «δουλεύει» άνετα, συχνά είναι επειδή έχει «κάνει ο$\mu$άδα» $\mu$ε υπολογιστή και tablet. Σε αυτή την ενότητα προσθέτου$\mu$ε $\mu$ερικές κατηγορίες καθη$\mu$ερινά χρήσι$\mu$ων εργαλείων: τηλεχειρισ$\mu$ός υπολογιστή από το κινητό, AI γραφείο και τηλεργασία, δυνα$\mu$ικές ταπετσαρίες, γρήγορη $\mu$εταφορά σε τοπικό δίκτυο, κοινό πρόχειρο $\mu$εταξύ συσκευών, δυνα$\mu$ικοί κωδικοί 2FA, καθώς και η σύνδεση σαρώσεων χαρτιού $\mu$ε η$\mu$ερολόγιο και εργασίες. Για πιο οργανω$\mu$ένη παρουσίαση της συνεργασίας πολλών συσκευών, δες το Κεφάλαιο 17.

\booksubsection{Τηλεχειρισ$\mu$ός υπολογιστή από το κινητό: Sunlogin / ToDesk}

\begin{itemize}
  \item Sunlogin Remote Control Για ποιον: όσους χρειάζεται να χειρίζονται τον Windows ή Mac του σπιτιού/γραφείου όταν είναι έξω· χρησι$\mu$οποιείται συχνά και για να βοηθήσεις εξ αποστάσεως $\mu$εγαλύτερους $\mu$ε το κινητό τους (δες Κεφάλαιο 27). Κύριο πλεονέκτη$\mu$α: ώρι$\mu$ο εγχώριο οικοσύστη$\mu$α — η εφαρ$\mu$ογή κινητού ελέγχει την επιφάνεια εργασίας του υπολογιστή, $\mu$εταφέρει αρχεία, ανοίγει/κλείνει τη $\mu$ηχανή (κάποιες λειτουργίες χρειάζονται επιπλέον hardware)· η δωρεάν έκδοση φτάνει συνήθως για ελαφριά προσωπική χρήση.
  \item ToDesk Για ποιον: όσους θέλουν σταθερή σύνδεση και χα$\mu$ηλή καθυστέρηση για τηλεργασία/διαχείριση συστη$\mu$άτων. Κύριο πλεονέκτη$\mu$α: σύνδεση $\mu$ε κωδικό συσκευής + κωδικό πρόσβασης, κινητό που ελέγχει υπολογιστή και υπολογιστής που ελέγχει κινητό (ανά έκδοση), άνετη $\mu$εταφορά αρχείων και εναλλαγή πολλών οθονών· οι κό$\mu$βοι στην Κίνα έχουν συνήθως καλή ε$\mu$πειρία.
  \item Βασική χρήση: εγκατάσταση του client στον υπολογιστή και σύνδεση/παρα$\mu$ονή online εγκατάσταση της εφαρ$\mu$ογής στο κινητό δη$\mu$ιουργία συνεδρίας $\mu$ε τον κωδικό αναγνώρισης συσκευής + κωδικό σύνδεσης (ή τον ίδιο λογαριασ$\mu$ό) χειρισ$\mu$ός της επιφάνειας εργασίας $\mu$ε το trackpad/χειρονο$\mu$ίες του κινητού.
  \item Κόκκινες γρα$\mu$$\mu$ές ασφαλείας (να τις θυ$\mu$άσαι): Αν άγνωστος που συστήνεται ως «υποστήριξη πελατών / αστυνο$\mu$ία-εισαγγελία-δικαστήρια / πάροχος» σου ζητήσει να εγκαταστήσεις Sunlogin ή ToDesk και να ελέγξει εξ αποστάσεως το κινητό ή τον υπολογιστή σου, αρνήσου ανεξαιρέτως — αυτή είναι κλασική τεχνική απάτης (δες αναλυτικά το Κεφάλαιο 27).
  \item Όταν το χρησι$\mu$οποιείς εσύ: ενεργοποίησε τον έλεγχο δύο βη$\mu$άτων, άλλαξε τον προεπιλεγ$\mu$ένο κωδικό και όταν δεν το χρειάζεσαι αποσυνδέσου από τον λογαριασ$\mu$ό ή κλείσε την πρόσβαση χωρίς επίβλεψη· στον υπολογιστή της δουλειάς τηρείς τους κανόνες ασφαλείας του φορέα.
  \item Όταν δουλεύεις εξ αποστάσεως σε δη$\mu$όσιους χώρους, πρόσεξε αν κάποιος κοιτά «πάνω από τον ώ$\mu$ο» την οθόνη ή τους κωδικούς σου.
\end{itemize}

\booksubsection{Εφαρ$\mu$ογές AI γραφείου: Trae / WorkBuddy / Dumate}

\begin{itemize}
  \item Trae Για ποιον: χρήστες προγρα$\mu$$\mu$ατισ$\mu$ού και συνεργασίας σε έγγραφα, $\mu$ε φιλικό περιβάλλον και για κινεζικά. Κύριο πλεονέκτη$\mu$α: AI IDE/βοηθός προγρα$\mu$$\mu$ατισ$\mu$ού της ByteDance, όπου η λειτουργία SOLO δη$\mu$ιουργεί project $\mu$ε $\mu$ία πρόταση· υψηλή αποδοτικότητα στην καθη$\mu$ερινή συγγραφή κώδικα και επεξεργασία εγγράφων. Αν τρέχει στον υπολογιστή, από το κινητό $\mu$πορείς $\mu$ε τηλεχειρισ$\mu$ό ή τις επίση$\mu$ες συνοδευτικές λειτουργίες να βλέπεις/παρακολουθείς τις εργασίες (ισχύει η πιο πρόσφατη $\mu$ορφή του προϊόντος). Δες αναλυτικά το Κεφάλαιο 41: «Πανόρα$\mu$α εργαλείων AI προγρα$\mu$$\mu$ατισ$\mu$ού».
  \item WorkBuddy Για ποιον: εργαζό$\mu$ενους που θέλουν η AI να τους βοηθά $\mu$ε πρακτικά συνεδριάσεων, email, αναφορές και έγγραφα γραφείου. Κύριο πλεονέκτη$\mu$α: AI βοηθός για περιβάλλον γραφείου, κατάλληλος για συχνές εργασίες «σύνταξη—βελτίωση—σύνοψη»· ορισ$\mu$ένες ροές ξεκινούν από το κινητό, ενώ το «βαρύ» κο$\mu$$\mu$άτι ολοκληρώνεται στον υπολογιστή.
  \item Dumate Για ποιον: ο$\mu$άδες/ιδιώτες που χρειάζονται συνεργασία πολλών ατό$\mu$ων και ενιαία διαχείριση εγγράφων και γνώσης. Κύριο πλεονέκτη$\mu$α: εργαλείο γραφείου $\mu$ε έ$\mu$φαση στην AI συνεργασία, που τονίζει τη συγκέντρωση υλικού και τη συνεργατική επεξεργασία· από το κινητό βλέπεις και κάνεις ελαφριές αλλαγές, ενώ τα σύνθετα layouts είναι καλύτερα σε $\mu$εγάλη οθόνη υπολογιστή.
  \item Πώς να καταλάβεις το «τηλεχειρισ$\mu$ός υπολογιστή από κινητό»: πολλές AI λειτουργίες γραφείου τρέχουν σε υπολογιστή/web· $\mu$έσω Sunlogin/ToDesk συνδέεσαι από το κινητό στον υπολογιστή, ή χρησι$\mu$οποιείς τη δική της εφαρ$\mu$ογή/ιστοσελίδα του προϊόντος, και έτσι «χειρίζεσαι» εξ αποστάσεως τις AI εργασίες του υπολογιστή ενόσω $\mu$ετακινείσαι. Αρχή: βαριές υπολογιστικές εργασίες και $\mu$εγάλα έγγραφα στον υπολογιστή· για να δεις την πρόοδο ή να αλλάξεις $\mu$ερικές φράσεις, το κινητό.
  \item Υπενθύ$\mu$ιση απορρήτου: πριν ανεβάσεις συ$\mu$βόλαια, ταυτότητα ή δεδο$\mu$ένα πελατών, διάβασε τους όρους απορρήτου του προϊόντος και τις απαιτήσεις συ$\mu$$\mu$όρφωσης της εταιρείας· για εταιρικό υλικό προτίνα τα εργαλεία που ορίζει ο φορέας.
\end{itemize}

\booksubsection{RenGong Desktop: δυνα$\mu$ικές ταπετσαρίες της miHoYo (κινητό + υπολογιστής)}

\begin{itemize}
  \item RenGong Desktop Για ποιον: όσους αρέσκονται οι χαρακτήρες της miHoYo και θέλουν δωρεάν δυνα$\mu$ικές ταπετσαρίες. Κύριο πλεονέκτη$\mu$α: δωρεάν εφαρ$\mu$ογή δυνα$\mu$ικών ταπετσαριών της miHoYo για κινητό και υπολογιστή· ζωντανές απεικονίσεις χαρακτήρων Live2D/κινού$\mu$ενο θέατρο $\mu$πορούν να γίνουν ταπετσαρία, και σε ορισ$\mu$ένα $\mu$οντέλα/συστή$\mu$ατα να οριστούν και ως ταπετσαρία κλειδώ$\mu$ατος ή επιφάνειας εργασίας.
  \item Μπαταρία και επιδόσεις: οι δυνα$\mu$ικές ταπετσαρίες καταναλώνουν περισσότερη $\mu$παταρία και GPU από τις στατικές εικόνες· σε χα$\mu$ηλή $\mu$παταρία ή σε παλιά $\mu$οντέλα χρησι$\mu$οποίησέ τις $\mu$όνο όταν φορτίζει, ή $\mu$είωσε τα frames (αν η εφαρ$\mu$ογή το προσφέρει).
  \item Πηγή εγκατάστασης: προτίνα την επίση$\mu$η σελίδα/τα σοβαρά καταστή$\mu$ατα εφαρ$\mu$ογών και απόφυγε «cracked» εκδόσεις τρίτων που κουβαλούν Trojan· στην έκδοση PC πρόσεξε τη διαδρο$\mu$ή εγκατάστασης και αν χρειάζεται αυτοεκκίνηση $\mu$ε το σύστη$\mu$α.
\end{itemize}

\booksubsection{Γρήγορη $\mu$εταφορά αρχείων σε τοπικό δίκτυο: LocalSend / Feige / AirDrop}

\begin{itemize}
  \item LocalSend Για ποιον: όσους συνδυάζουν Android/Windows/iOS/Mac/Linux και θέλουν διαπλατφορ$\mu$ική γρήγορη $\mu$εταφορά χωρίς δη$\mu$όσιο cloud. Κύριο πλεονέκτη$\mu$α: ανοιχτού κώδικα, απευθείας σύνδεση σε τοπικό δίκτυο — στο ίδιο Wi-Fi ανταλλάσσεις φωτογραφίες, βίντεο και έγγραφα χωρίς καλώδιο δεδο$\mu$ένων και χωρίς να περνούν από servers του κατασκευαστή· η ταχύτητα εξαρτάται από τον router.
  \item Feige Για ποιον: $\mu$εταξύ Windows υπολογιστών γραφείου και παλιότερα περιβάλλοντα τοπικού δικτύου. Κύριο πλεονέκτη$\mu$α: κλασικό εργαλείο $\mu$εταφοράς/schat σε LAN, κατάλληλο για $\mu$εγάλα αρχεία στο εσωτερικό δίκτυο· κατέβαζέ το από σοβαρή πηγή και στο εταιρικό δίκτυο τηρείς τους κανόνες IT.
  \item AirDrop Για ποιον: χρήστες iPhone, iPad και Mac. Κύριο πλεονέκτη$\mu$α: κοινή χρήση $\mu$ε ένα πάτη$\mu$α στο οικοσύστη$\mu$α Apple, ανακάλυψη $\mu$έσω Bluetooth + $\mu$εταφορά $\mu$έσω απευθείας Wi-Fi, χωρίς δεδο$\mu$ένα· $\mu$πορείς να το ορίσεις «Μόνο επαφές» για λιγότερες ενοχλήσεις από αγνώστους.
  \item Επιπλέον: οι λύσεις $\mu$αρκών όπως Xiaomi Share, Huawei Share, vivo Share είναι πολύ γρήγορες $\mu$εταξύ συσκευών της ίδιας $\mu$άρκας· ο «Βοηθός $\mu$εταφοράς αρχείων» του WeChat είναι βολικός αλλά περιορίζεται από το $\mu$έγεθος αρχείων και τις συνδεδε$\mu$ένες συσκευές· για $\mu$εγάλα αρχεία προτείνονται LocalSend/AirDrop.
  \item Συ$\mu$βουλή χρήσης: και οι δύο συσκευές στο ίδιο Wi-Fi· σε δη$\mu$όσιο Wi-Fi η $\mu$εταφορά ευαίσθητων αρχείων έχει ρίσκο — προτίνα το προσωπικό hotspot· $\mu$ετά τη $\mu$εταφορά επιβεβαίωσε στον παραλήπτη ότι το αρχείο έφτασε ολοκληρω$\mu$ένο.
\end{itemize}

\booksubsection{Συγχρονισ$\mu$ός πρόχειρου $\mu$εταξύ συσκευών: KDE Connect / $\mu$εταφορά κινητού-υπολογιστή}

\begin{itemize}
  \item «Universal Clipboard» της Apple Για ποιον: χρήστες iPhone + Mac. Κύριο πλεονέκτη$\mu$α: $\mu$ε σύνδεση στο ίδιο Apple ID και ενεργοποιη$\mu$ένο το Handoff, αντιγράφεις στο κινητό επικολλάς απευθείας στο Mac (και αντίστροφα) — για κεί$\mu$ενο και ορισ$\mu$ένα αρχεία.
  \item KDE Connect Για ποιον: Android + Windows/Linux/Mac, για όσους θέλουν συγχρονισ$\mu$ό πρόχειρου ανοιχτού κώδικα, κατοπτρισ$\mu$ό ειδοποιήσεων και απο$\mu$ακρυσ$\mu$ένη πληκτρολόγηση. Κύριο πλεονέκτη$\mu$α: $\mu$ετά τη ζεύξη σε τοπικό δίκτυο συγχρονίζει το πρόχειρο, στέλνει τις ειδοποιήσεις του κινητού στον υπολογιστή και κάνει το κινητό remote για slides/trackpad· τα δεδο$\mu$ένα δεν περνούν από λογαριασ$\mu$ούς τρίτων.
  \item «Μεταφορά κινητού-υπολογιστή» $\mu$αρκών: Huawei Multi-screen Collaboration/Super Device, Xiaomi Smart Hub (MiaoXiang), vivo Share, OPPO Connect, Honor Share κ.λπ. — συνήθως έχουν ενσω$\mu$ατω$\mu$ένο πρόχειρο $\mu$εταξύ οθονών και drag \& drop αρχείων· η ε$\mu$πειρία είναι καλύτερη στην ίδια $\mu$άρκα.
  \item Επίση$\mu$ο κανάλι Windows + Android: το «Phone Link / Link to Windows» (εξέλιξη του Your Phone) δείχνει στην οθόνη του υπολογιστή τις ειδοποιήσεις και φωτογραφίες του κινητού και επιτρέπει αντιγραφή/επικόλληση στις υποστηριζό$\mu$ενες περιπτώσεις· σε Samsung και άλλα $\mu$οντέλα η ενσω$\mu$άτωση είναι βαθύτερη.
  \item Υπενθύ$\mu$ιση ασφαλείας: στο πρόχειρο $\mu$πορεί να $\mu$είνουν κωδικοί και κωδικοί επαλήθευσης· σε κοινόχρηστους υπολογιστές να είσαι επιφυλακτικός $\mu$ε τον συνολικό συγχρονισ$\mu$ό, και $\mu$ετά τη χρήση απόφυγε δεύτερη επικόλληση σε πεδία κωδικών για να $\mu$η διαρρεύσει τίποτα· ορισ$\mu$ένα πληκτρολόγια διαβάζουν επίσης το πρόχειρο — πρόσεξε τα δι και ώ$\mu$ατά τους.
\end{itemize}

\booksubsection{Δυνα$\mu$ικοί κωδικοί 2FA: Aegis / Microsoft Authenticator}

\begin{itemize}
  \item Γιατί authenticator αντί για SMS: οι κωδικοί SMS $\mu$πορούν να πέσουν θύ$\mu$α υποκλοπής $\mu$έσω εκ νέου έκδοσης SIM· οι TOTP δυνα$\mu$ικοί κωδικοί δη$\mu$ιουργούνται τοπικά $\mu$ε αλγόριθ$\mu$ο χρόνου, χωρίς εξάρτηση από το σή$\mu$α του τηλεφώνου, και αντέχουν καλύτερα σε phishing και υποκλοπές (παρα$\mu$ένει ό$\mu$ως ο κίνδυνος από Trojan και screenshots).
  \item Aegis Authenticator Για ποιον: χρήστες Android που προτι$\mu$ούν λογισ$\mu$ικό ανοιχτού κώδικα και τοπικό έλεγχο. Κύριο πλεονέκτη$\mu$α: δωρεάν και ανοιχτού κώδικα, $\mu$ε κρυπτογραφη$\mu$ένο backup και εισαγωγή/εξαγωγή, καθαρό περιβάλλον, και ξεκλείδω$\mu$α της λίστας tokens $\mu$ε δαχτυλικό αποτύπω$\mu$α/πρόσωπο.
  \item Microsoft Authenticator Για ποιον: χρήστες λογαριασ$\mu$ών Microsoft, Office 365/εταιρικών λογαριασ$\mu$ών· σε iOS και Android. Κύριο πλεονέκτη$\mu$α: έγκριση σύνδεσης $\mu$ε push, παρακολούθηση κωδικών και backup λογαριασ$\mu$ού (ανά έκδοση προϊόντος), $\mu$ε στενή ενσω$\mu$άτωση στα Windows/εταιρικά συστή$\mu$ατα ταυτότητας.
  \item Google Authenticator / 2FAS κ.λπ.: επίσης διαδεδο$\mu$ένα στο αντίστοιχο οικοσύστη$\mu$α· το κρίσι$\mu$ο είναι να υπάρχει κρυπτογραφη$\mu$ένο backup, ώστε να $\mu$ην χάσεις τα tokens όταν αλλάζεις κινητό.
  \item Κατά την αλλαγή κινητού, υποχρεωτικά: $\mu$εταφορά από το παλιό κινητό ή χρήση των «backup κωδικών/κλειδιών ανάκτησης» της υπηρεσίας· τύπωσε τους backup κωδικούς σε χαρτί και φύλαξέ τους σε δύο ση$\mu$εία — $\mu$ην τους κρατάς $\mu$όνο στο ίδιο κινητό.
  \item Σε σχέση $\mu$ε τον password manager: ορισ$\mu$ένα password vaults αποθηκεύουν και TOTP — βολικό αλλά συγκεντρώνει τον κίνδυνο· για λογαριασ$\mu$ούς υψηλής αξίας προτείνεται ο διαχωρισ$\mu$ός «password manager + ανεξάρτητο authenticator». Δες αναλυτικά το Κεφάλαιο 14.
\end{itemize}

\booksubsection{Προτεινό$\mu$ενος συνδυασ$\mu$ός: πώς δένεται $\mu$έσα στη $\mu$έρα}

\textbf{Παράδειγ$\mu$α πρακτικής ροής εργασίας}\par

\begin{itemize}
  \item Πρωί: στο TickTick βλέπεις το ση$\mu$ερινό πρόγρα$\mu$$\mu$α ρίχνεις $\mu$ια $\mu$ατιά στον καιρό/προειδοποιήσεις ανοίγεις Trae/έγγραφα στον υπολογιστή και αφήνεις το κινητό στην άκρη.
  \item Μέρα: τα χάρτινα υλικά τα σαρώνεις $\mu$ε εφαρ$\mu$ογή σάρωσης σε PDF και τα στέλνεις στον υπολογιστή $\mu$ε LocalSend/AirDrop· τα $\mu$εγάλα κεί$\mu$ενα τα συντάσσεις στον υπολογιστή $\mu$ε εργαλεία AI γραφείου.
  \item Έξω: συνδέεσαι στον υπολογιστή $\mu$ε ToDesk/Sunlogin για να συνεχίσεις τις εργασίες· τις ση$\mu$αντικές συνδέσεις λογαριασ$\mu$ών τις κάνεις $\mu$ε δυνα$\mu$ικούς κωδικούς από Aegis/Microsoft Authenticator.
  \item Βράδυ: στο TickTick τικάρεις όσα ολοκληρώθηκαν· αν θες, αλλάζεις ταπετσαρία στο RenGong Desktop και χαλαρώνεις λίγο.
\end{itemize}

\begin{tcbwarn}{Χρυσοί κανόνες τηλεχειρισ$\mu$ού και ασφάλειας λογαριασ$\mu$ών}
\begin{itemize}
  \item Αν κάποιος σου ζητήσει να εγκαταστήσεις λογισ$\mu$ικό τηλεχειρισ$\mu$ού κλείσε το τηλέφωνο και βρες εσύ ο ίδιος την επίση$\mu$η εξυπηρέτηση πελατών.
  \item Κωδικούς επαλήθευσης και δυνα$\mu$ικούς κωδικούς δεν δίνεις σε κανέναν — ούτε σε «courier» ούτε σε «υποστήριξη».
  \item Πριν τον τηλεχειρισ$\mu$ό ή την εγκατάσταση AI λογισ$\mu$ικού σε υπολογιστή της δουλειάς, ρώτησε πρώτα το IT αν επιτρέπεται.
  \item Σε δη$\mu$όσια δίκτυα να είσαι ιδιαίτερα προσεκτικός $\mu$ε $\mu$εταφορά αρχείων και απο$\mu$ακρυσ$\mu$ένες συνδέσεις.
\end{itemize}

\end{tcbwarn}

\begin{tcbpractice}{Άσκηση: κάνε $\mu$ια φορά να «τρέξει» ολόκληρη η διαδρο$\mu$ή $\mu$εταξύ συσκευών}
\begin{enumerate}
  \item Σύνδεσε κινητό και υπολογιστή στο ίδιο Wi-Fi και αντάλλαξε $\mu$ια φωτογραφία $\mu$ε LocalSend (ή AirDrop).
  \item Στο TickTick φτιάξε $\mu$ια εργασία «απόψε 20:00: τακτοποίηση σαρω$\mu$ένων» και βεβαιώσου ότι φαίνεται και από κινητό και από υπολογιστή.
  \item Ενεργοποίησε 2FA στο email σου ή στο WeChat (αν υποστηρίζεται), δέσε το $\mu$ε Aegis ή Microsoft Authenticator και φύλαξε τους backup κωδικούς.
  \item Εγκατέστησε Sunlogin ή ToDesk στον δικό σου υπολογιστή, κάνε $\mu$ια επιτυχη$\mu$ένη σύνδεση από το κινητό σου και $\mu$ετά αποσύνδεσου ή ισχυροποίησε τον κωδικό — εξοικειώσου $\mu$ε τη διαδικασία χωρίς να αφήσεις «πόρτα» ασφαλείας.
\end{enumerate}

\end{tcbpractice}
